\documentclass[11pt]{article}
\usepackage{amsmath,amssymb}
\usepackage{graphicx}
\usepackage{hyperref}
\usepackage[numbers]{natbib}

% Macros used throughout the paper.
\newcommand{\G}{\mathcal{G}}
\newcommand{\E}{\mathbb{E}}

\title{In-Context Learning in Graph Neural Networks}
\author{Anonymous Authors}

\begin{document}
\maketitle

\begin{abstract}
Large language models can learn new tasks from a handful of examples supplied in their prompt, a capability known as in-context learning. Whether graph neural networks admit a comparable mechanism remains unclear. We study in-context learning for node classification on graphs $\G = (V, E)$ and show that a simple prompting scheme lets a frozen model adapt to unseen label sets without any gradient updates.
\end{abstract}

\section{Introduction}\label{sec:intro}
Graph neural networks (GNNs) are the standard tool for learning on relational data \cite{kipf2017gcn, velickovic2018gat}. Their success, however, rests on task-specific training: a model trained to predict one set of labels must be fine-tuned before it can predict another.

In contrast, large language models perform new tasks when shown a few input-output pairs in context \cite{brown2020gpt3}. This raises a natural question: can a graph model be prompted in the same way?

We answer this question affirmatively. Our contributions are threefold. First, we formalize in-context learning for graphs. Second, we propose a prompt graph construction that connects query nodes to labelled example nodes. Third, we show empirically that pretrained GNNs equipped with prompt graphs match fine-tuned baselines on four benchmarks (Section~\ref{sec:experiments}).

\section{Related Work}
\paragraph{Prompting for graphs.} Several recent works adapt prompting ideas to graph models, typically by learning soft prompt vectors that are prepended to node features. These methods still require optimization for each new task.

\paragraph{Meta-learning.} Few-shot node classification has also been approached with meta-learning, where a model is trained over many sampled tasks so that it can adapt quickly. Unlike our approach, these methods update parameters at test time.

\section{Method}\label{sec:method}
Let $X \in \mathbb{R}^{n \times d}$ denote node features and $A$ the adjacency matrix. A message-passing layer computes
\begin{equation}
  h_v^{(\ell+1)} = \sigma\Big( W^{(\ell)} \sum_{u \in \mathcal{N}(v)} \frac{h_u^{(\ell)}}{\sqrt{d_u d_v}} \Big),
  \label{eq:mp}
\end{equation}
where $\mathcal{N}(v)$ is the neighbourhood of $v$ and $\sigma$ is a nonlinearity.

Given $k$ labelled examples per class, we build a prompt graph by adding one virtual node per class and connecting each example to the virtual node of its label. A query node is classified by comparing its representation with those of the virtual nodes, as in Eq.~\eqref{eq:mp}.

\begin{figure}[t]
  \centering
  \includegraphics[width=0.8\linewidth]{figures/prompt_graph.pdf}
  \caption{Construction of the prompt graph. Example nodes are linked to virtual label nodes; the query node attends to all of them.}
  \label{fig:prompt}
\end{figure}

\subsection{Training Objective}
We pretrain the model with a contrastive objective that pulls each node towards the virtual node of its class and pushes it away from the others:
\begin{align}
  \mathcal{L} &= -\E_{v}\left[ \log \frac{\exp(s(v, c_v)/\tau)}{\sum_{c} \exp(s(v, c)/\tau)} \right].
\end{align}
The temperature $\tau$ is fixed to $0.1$ in all experiments.

\section{Experiments}\label{sec:experiments}
We evaluate on Cora, CiteSeer, PubMed and ogbn-arxiv. For each dataset, we hold out a disjoint set of classes that is never seen during pretraining.

\begin{table}[t]
  \centering
  \caption{Few-shot node classification accuracy (\%) on held-out classes.}
  \begin{tabular}{lcccc}
    Method & Cora & CiteSeer & PubMed & arxiv \\
    Fine-tuning & 71.2 & 63.5 & 70.1 & 55.4 \\
    Ours (in-context) & 70.8 & 64.1 & 69.7 & 56.0 \\
  \end{tabular}
  \label{tab:main}
\end{table}

Table~\ref{tab:main} shows that in-context prediction is competitive with fine-tuning while requiring no gradient steps. The gap is largest on PubMed, where only three classes exist and the prompt graph is correspondingly small.

% TODO: add ablation on the number of examples per class.

\section{Conclusion}
We showed that graph neural networks can learn in context when the task is expressed as a prompt graph. Future work will study larger pretrained models and heterogeneous graphs.

\bibliographystyle{plainnat}
\bibliography{refs}

\end{document}
